\chapter{Conclusão}
\label{sec:conclusao}

Este trabalho apresentou o design e a implementação de uma plataforma de processamento de mídia genérica, cuja concepção foi guiada pela necessidade de superar as limitações de sistemas tradicionais, como o alto acoplamento e a baixa extensibilidade. O problema foi abordado em duas frentes: a do usuário final, que enfrenta a sobrecarga de informação, e a do engenheiro de software, que enfrenta os desafios de construir sistemas complexos, manuteníveis e escaláveis.

A solução foi materializada como um monólito modular, orientado a eventos, que implementa um pipeline de processamento desacoplado. A abstração do conceito de "vídeo" para uma entidade de "Mídia" genérica, combinada com a comunicação assíncrona via eventos, permitiu a criação de um sistema flexível, capaz de processar diferentes formatos de entrada (vídeo, áudio, texto) através de processadores especializados e independentes.

Os objetivos específicos, que representavam os pilares arquiteturais do projeto, foram alcançados com sucesso. A aplicação rigorosa do padrão \textbf{CQRS} em todos os módulos de negócio resultou em um código mais limpo e focado. O uso do \textbf{Padrão State} garantiu a integridade do ciclo de vida da entidade \texttt{Media}, enquanto os padrões \textbf{Strategy} e \textbf{Factory} permitiram a criação de um sistema de provedores dinâmicos, conferindo alta flexibilidade tecnológica à plataforma. A resiliência foi endereçada com o \textbf{Circuit Breaker Pattern}, e a observabilidade completa foi alcançada através da integração de logs estruturados, métricas Prometheus e tracing distribuído com OpenTelemetry, o que se provou essencial para a depuração de um sistema assíncrono.

A principal contribuição deste trabalho reside na síntese e aplicação prática de um conjunto de padrões de design e arquitetura de software modernos para resolver um problema do mundo real. O projeto demonstra que, ao priorizar o baixo acoplamento, a alta coesão e a separação de responsabilidades desde o início, é possível construir sistemas complexos que são, ao mesmo tempo, robustos, escaláveis e, crucialmente, fáceis de manter e evoluir.

Como trabalhos futuros, vislumbra-se a expansão do pipeline com novos processadores, como um serviço de análise de sentimento sobre o texto transcrito ou um serviço de reconhecimento de objetos em vídeos. A arquitetura atual, por ser modular e orientada a eventos, está perfeitamente preparada para acomodar essas novas funcionalidades com o mínimo de atrito. Além disso, a estratégia de testes pode ser expandida para incluir testes de carga e performance, validando a escalabilidade do sistema sob alta demanda. Finalmente, a infraestrutura de implantação poderia evoluir de Docker Compose para uma plataforma de orquestração mais robusta, como Kubernetes, para a qual a arquitetura baseada em contêineres já está preparada.
